% Bài làm buổi 17/11
\begin{problem}[9DS-HSB2-010]
a) Thay $v = 2$, $F = 120$ vào $F = av^2$: $120 = a \cdot 2^2 = 4a$, suy ra $a = 30$.

b) Với $v = 10$: $F = 30 \cdot 10^2 = 3\,000$ (N).

c) Đổi đơn vị: $90$ km/h $= \dfrac{90\,000}{3\,600}$ m/s $= 25$ m/s.

Lực của gió khi đó: $F = 30 \cdot 25^2 = 30 \cdot 625 = 18\,750$ (N).

Vì $18\,750 > 12\,000$ nên cánh buồm không chịu được, thuyền không thể đi trong gió bão $90$ km/h.

Đáp số: a) $a = 30$; b) $F = 3\,000$ N; c) Không, vì $F = 18\,750$ N $> 12\,000$ N
\end{problem}

\begin{problem}[9DS-HSB2-006]
a) $M(2; -2)$ thuộc $(P)$ nên $-2 = a \cdot 2^2 = 4a$, suy ra $a = -\dfrac{1}{2}$. Vậy $(P)\colon y = -\dfrac{1}{2}x^2$.

Vì $a = -\dfrac{1}{2} < 0$ nên $(P)$ nằm phía dưới trục hoành.

b) Phương trình hoành độ giao điểm: $-\dfrac{1}{2}x^2 = x - 4$, tức là $x^2 + 2x - 8 = 0$.

$\Delta' = 1^2 - 1 \cdot (-8) = 9$, $\sqrt{\Delta'} = 3$, nên $x_1 = -1 + 3 = 2$ và $x_2 = -1 - 3 = -4$.

Với $x = 2$: $y = -\dfrac{1}{2} \cdot 2^2 = -2$. Với $x = -4$: $y = -\dfrac{1}{2} \cdot (-4)^2 = -\dfrac{1}{2} \cdot (-16) = 8$.

Đáp số: a) $a = -\dfrac{1}{2}$, $(P)$ nằm phía dưới trục hoành; b) $(2; -2)$ và $(-4; 8)$
\end{problem}
